\documentclass[11pt,a4paper]{article}
\usepackage[utf8]{inputenc}
\usepackage[margin=1in]{geometry}
\usepackage{amsmath,amssymb,amsthm}
\usepackage{booktabs}
\usepackage{url}
\usepackage{color}
\usepackage{enumitem}

\newcommand{\coloneqq}{\mathrel{\mathop:}=}

\newtheorem{theorem}{Theorem}
\newtheorem{lemma}[theorem]{Lemma}
\newtheorem{proposition}[theorem]{Proposition}
\newtheorem{corollary}[theorem]{Corollary}
\newtheorem{definition}{Definition}
\newtheorem{remark}{Remark}

\title{\textbf{On the Asymptotic Validity of Gaussian Mixture Confidence Sequences for Model Cascade Calibration: Martingale Limit Theory, Diagnostic Failure Analysis, and Principled Safeguards}}
\author{ACTIS Technical Research Group}
\date{September 2026}

\begin{document}
\maketitle

\begin{abstract}
Asymptotic Gaussian mixture test supermartingales provide a powerful, closed-form framework for sequential hypothesis testing and anytime-valid inference. In two-threshold model cascade calibration, they enable substantial reductions in expensive ground-truth oracle queries compared to conservative finite-sample betting supermartingales. However, when deployed in conjunction with importance sampling under high recall constraints ($\gamma_R = 0.95$), recent empirical benchmarks revealed severe efficiency collapse: importance sampling variants requested excessive oracle queries, frequently being outperformed by uniform sampling without replacement and finite-sample baselines. 

In this research note, we present a rigorous theoretical investigation into the foundations of asymptotic confidence sequences in adaptive model cascades. First, we provide a complete formalization of the exact multi-threshold family of test martingales and identify the active certified threshold pair evaluated at stopping. Second, we revisit the Central Limit Theorem from first principles, establishing that sequential stopping in cascades is governed strictly by the \emph{Martingale Central Limit Theorem} (Hall and Heyde, 1980; McLeish, 1974) rather than classical independent and identically distributed limit theorems. We clarify the essential conditions required for asymptotic validity: predictable quadratic variation stabilization, the conditional Lindeberg condition, and jump dispersion bounds. Third, we critically evaluate Kish's Effective Sample Size (ESS) in the context of martingales, demonstrating that applying Kish's formula to sequential test martingales constitutes a double-counting of weight dispersion because the realized quadratic variation accumulator $V_t$ already accounts for importance weights exactly. Fourth, we deconstruct the exact mathematical mechanism behind the observed empirical degradation: an artificial dimensional scale mismatch between the sample variance accumulator $V_t$ and the zero-variance floor $v_{0, \text{floor}}$, which diverges by a factor of $((1 - \gamma_R)/\gamma_R)^2 \approx 1/361$ on accurate candidate thresholds. Finally, we formulate a principled, unified diagnostic framework based on expected cumulative process variance $v_{0, \text{target}}$ and Hall--Heyde jump dispersion ratios, eliminating ad-hoc piecewise heuristics. Extensive empirical evaluations across all four benchmarks in the SUPG suite confirm that this principled diagnostic completely restores the superior efficiency of importance sampling while strictly preserving target error guarantees ($0.0\%$ empirical violations across all benchmarks).
\end{abstract}

\tableofcontents
\newpage

\section{Introduction and Problem Setting}

Modern machine learning systems increasingly utilize multi-stage model cascades to optimize the trade-off between inference accuracy and operational cost. In a typical two-threshold cascade, a computationally cheap proxy model score $s(x) \in [0, 1]$ is evaluated on incoming items. Items exhibiting high proxy confidence ($s(x) \ge \tau_{\text{pos}}$) are automatically accepted as positive without human review; items exhibiting very low proxy confidence ($s(x) < \tau_{\text{neg}}$) are automatically rejected as negative; and ambiguous items falling in the intermediate routing band $[\tau_{\text{neg}}, \tau_{\text{pos}})$ are deferred to an expensive ground-truth oracle $y(x) \in \{0, 1\}$.

Calibrating such cascades requires selecting a threshold pair $(\tau_{\text{pos}}, \tau_{\text{neg}})$ over a finite transductive population $\mathcal{D} = \{(x_i, y_i)\}_{i=1}^N$ that minimizes oracle consultations while guaranteeing, with high confidence $1 - \delta$, that the cascade simultaneously satisfies pre-specified target recall $\gamma_R$ and precision $\gamma_P$:
\begin{equation}\label{eq:joint_guarantee}
\mathbb{P}\Big( \text{Recall}(\tau_{\text{neg}}) \ge \gamma_R \;\land\; \text{Precision}(\tau_{\text{pos}}, \tau_{\text{neg}}) \ge \gamma_P \Big) \ge 1 - \delta.
\end{equation}

\subsection{Asymptotic vs.\ Finite-Sample Confidence Sequences}
Sequential calibration methods draw samples $Z_1, Z_2, \dots$ adaptively and track non-negative test supermartingales $M_t$ to construct anytime-valid confidence sequences. Under Ville's inequality~\cite{ville1939etude}, the probability that $M_t$ ever crosses the significance threshold $1/\alpha$ under the null hypothesis is bounded by $\alpha$:
\begin{equation}\label{eq:ville_ineq}
\mathbb{P}_{H_0}\big(\exists t \ge 1 : M_t \ge 1/\alpha\big) \le \alpha.
\end{equation}

Two primary paradigms exist for constructing $M_t$:
\begin{enumerate}[leftmargin=*]
    \item \textbf{Finite-Sample Betting Supermartingales}: Formulated by Waudby-Smith and Ramdas~\cite{waudby2024}, these supermartingales require increments to be strictly bounded, typically scaled into $[0, 1]$. While non-asymptotically exact, this scaling introduces severe conservatism. When importance sampling is used to oversample rare positive items, importance weights $w(Z_t) = \frac{1}{N q(Z_t)}$ vary across multiple orders of magnitude (e.g., $w(x) \in [0.12, 2.4]$). Normalizing the increments by the global maximum weight $w_{\max}$ shrinks the true positive signal to near zero (e.g., $(1 - \gamma_R) / w_{\max} \approx 0.0038$). Consequently, the betting capital fails to compound, preventing early stopping and causing catastrophic efficiency collapse (querying nearly $100\%$ of the population).
    \item \textbf{Asymptotic Gaussian Mixture Supermartingales}: Derived from Robbins~\cite{robbins1970} and Howard et al.~\cite{howard2021}, mixture supermartingales integrate likelihood ratios against a Gaussian prior $\mathcal{N}(0, v_0)$. They operate directly on the unnormalized sample sum $S_t = \sum_{i=1}^t X_i$ and the empirical variance accumulator $V_t = \sum_{i=1}^t (X_i - \bar{X}_t)^2$, offering closed-form algebraic invertibility and high power. However, their coverage guarantee is asymptotic, relying on normal approximations of the boundary crossing process.
\end{enumerate}

\subsection{The Empirical Performance Anomaly}
To combine the variance-reduction benefits of importance sampling with the superior power of Gaussian mixture confidence sequences, the ACTIS framework implements an asymptotic importance-sampling calibration method (\texttt{actis\_adaptive\_is\_asymptotic}). 

However, recent evaluations on standard SUPG benchmarks~\cite{kang2020approximate} revealed an unexpected anomaly: at confidence levels $\delta = 0.20$ and $\delta = 0.05$, the importance sampling variant exhibited worse oracle query efficiency than uniform sampling without replacement (\texttt{actis\_adaptive}) and even lagged behind the finite-sample baseline BARGAIN-PR on several datasets. Specifically:
\begin{itemize}[leftmargin=*]
    \item On \textbf{SUPG OntoNotes} ($\delta = 0.20$), \texttt{actis\_adaptive\_is\_asymptotic} required an oracle call rate of $47.56\%$, whereas BARGAIN-PR required $42.93\%$.
    \item On \textbf{SUPG TACRED} ($\delta = 0.05$), \texttt{actis\_adaptive\_is\_asymptotic} consumed $49.32\%$ oracle calls, whereas uniform ACTIS required only $40.14\%$.
\end{itemize}
Detailed profiling traces demonstrated that the sequential value-of-information stopping criterion was repeatedly overridden by a runtime diagnostic component, \path{AsymptoticValidityDiagnostics}. The diagnostic consistently flagged $V_t(R) / v_0 < 0.05$, declaring \texttt{asymp\_valid = False} and forcing the algorithm to continue drawing calibration queries until reaching the maximum allocation limit $\text{max\_sample\_size} = 0.50 N$.

\subsection{Contributions of this Research Note}
In response to these findings, this note provides:
\begin{enumerate}[leftmargin=*]
    \item A complete mathematical specification of the underlying families of test martingales across the threshold search lattice, clarifying precisely which martingales are certified at stopping.
    \item A rigorous grounding of asymptotic anytime validity in the \emph{Martingale Central Limit Theorem} (Hall and Heyde, 1980; McLeish, 1974), detailing the exact necessary conditions: predictable quadratic variation consistency, the conditional Lindeberg condition, and jump dispersion.
    \item A critical analysis of Kish's Effective Sample Size, demonstrating why Kish's formula applies strictly to static survey sampling and does not apply to sequential martingales where realized quadratic variation $V_t$ already incorporates squared weights.
    \item An analytical proof that the diagnostic failure was caused by an dimensional scale mismatch between $V_t$ and the zero-variance safety floor $v_{0, \text{floor}}$, which diverges by a factor of $((1 - \gamma_R)/\gamma_R)^2 \approx 1/361$ when candidate thresholds exhibit zero empirical false negatives ($n_{\text{FN}} = 0$).
    \item A principled diagnostic framework based on expected process variance $v_{0, \text{target}}$ and Hall--Heyde jump dispersion bounds that eliminates ad-hoc piecewise branching, fully restoring the empirical superiority of importance sampling across all SUPG benchmarks.
\end{enumerate}

\newpage
\section{Formalization of the Test Martingales}

We consider a transductive population $\mathcal{D} = \{(x_i, y_i)\}_{i=1}^N$. Let $s(x) \in [0, 1]$ be the proxy model score. Candidate lower thresholds are chosen from a discretized grid $\mathcal{T}_l = \{\tau_0, \tau_1, \dots, \tau_{K_l-1}\}$ with $\tau_0 = 0.0$, and candidate upper thresholds from $\mathcal{T}_u = \{\tau_0', \tau_1', \dots, \tau_{K_u-1}'\}$ with $\tau_{K_u-1}' = 1.0$.

Items $Z_1, Z_2, \dots$ are sampled independently with replacement from a proposal distribution $q(x) > 0$. The sequential filtration is defined by $\mathcal{F}_t = \sigma(Z_1, Y_1, \dots, Z_t, Y_t)$ with $\mathcal{F}_0$ containing all prior population quantities (including proxy scores and proposal weights). Unbiased importance weights are given by $w(x) = \frac{1}{N q(x)}$.

\subsection{The Recall Martingale Family}
For each candidate lower threshold index $l \in \{0, \dots, K_l-1\}$, the linearized population recall margin is:
\begin{equation}\label{eq:recall_margin}
\mu_R(l) \coloneqq \frac{1}{N}\sum_{i=1}^N Y_i \big(\mathbf{1}\{s(x_i) \ge \tau_l\} - \gamma_R\big).
\end{equation}
Testing $\text{Recall}(\tau_l) \ge \gamma_R$ is mathematically equivalent to testing the one-sided null hypothesis:
\begin{equation}\label{eq:h0_recall}
H_0^{(R)}(l) : \mu_R(l) \le 0 \quad \text{versus} \quad H_1^{(R)}(l) : \mu_R(l) > 0.
\end{equation}

For each candidate lower threshold $\tau_l$, the sequence of observations is:
\begin{equation}\label{eq:obs_recall}
X_{t, l}^{(R)} \coloneqq w(Z_t) Y_t \big(\mathbf{1}\{s(Z_t) \ge \tau_l\} - \gamma_R\big).
\end{equation}
By importance sampling construction, $\mathbb{E}[X_{t, l}^{(R)} \mid \mathcal{F}_{t-1}] = \mu_R(l)$. Under the boundary null $\mu_R(l) = 0$, $X_{t, l}^{(R)}$ is a zero-mean martingale difference sequence.

We define the running sum and sample variance accumulator for candidate threshold $l$:
\begin{align}
S_t^{(R)}(l) &\coloneqq \sum_{i=1}^t X_{i, l}^{(R)}, \label{eq:st_recall} \\
V_t^{(R)}(l) &\coloneqq \sum_{i=1}^t \big(X_{i, l}^{(R)} - \bar{X}_{t, l}^{(R)}\big)^2 \quad \text{where } \bar{X}_{t, l}^{(R)} = \frac{1}{t} S_t^{(R)}(l). \label{eq:vt_recall}
\end{align}

Let $v_{0, R}[l]$ denote the mixture prior variance for threshold $l$. The Gaussian mixture test supermartingale for recall at threshold $l$ is given in closed form by:
\begin{equation}\label{eq:gm_recall}
M_t^{(R)}(l) \coloneqq \sqrt{\frac{v_{0, R}[l]}{v_{0, R}[l] + V_t^{(R)}(l)}} \exp\left( \frac{\big(\max\{0, S_t^{(R)}(l)\}\big)^2}{2\big(v_{0, R}[l] + V_t^{(R)}(l)\big)} \right).
\end{equation}
Rejection of $H_0^{(R)}(l)$ occurs at the first stopping time $\tau$ where:
\begin{equation}\label{eq:reject_recall}
M_t^{(R)}(l) \ge \frac{1}{\delta_R} \quad \text{with } \delta_R = \frac{\delta}{2}.
\end{equation}

\subsection{The Precision Martingale Family}
For each candidate threshold pair $(u, l)$ with $0 \le l \le u < K_u$, the linearized population precision margin is:
\begin{equation}\label{eq:precision_margin}
\mu_P(u, l) \coloneqq \frac{1}{N}\sum_{i=1}^N \Big( (1 - \gamma_P) Y_i \mathbf{1}\{s(x_i) \ge \tau_l\} - \gamma_P (1 - Y_i) \mathbf{1}\{s(x_i) \ge \tau_u'\} \Big).
\end{equation}
Testing $\text{Precision}(\tau_u', \tau_l) \ge \gamma_P$ corresponds to testing:
\begin{equation}\label{eq:h0_precision}
H_0^{(P)}(u, l) : \mu_P(u, l) \le 0 \quad \text{versus} \quad H_1^{(P)}(u, l) : \mu_P(u, l) > 0.
\end{equation}

The sequential observation increment is:
\begin{equation}\label{eq:obs_precision}
X_{t, u, l}^{(P)} \coloneqq w(Z_t) \Big( (1 - \gamma_P) Y_t \mathbf{1}\{s(Z_t) \ge \tau_l\} - \gamma_P (1 - Y_t) \mathbf{1}\{s(Z_t) \ge \tau_u'\} \Big).
\end{equation}
The running sum and sample variance accumulator are:
\begin{align}
S_t^{(P)}(u, l) &\coloneqq \sum_{i=1}^t X_{i, u, l}^{(P)}, \label{eq:st_precision} \\
V_t^{(P)}(u, l) &\coloneqq \sum_{i=1}^t \big(X_{i, u, l}^{(P)} - \bar{X}_{t, u, l}^{(P)}\big)^2. \label{eq:vt_precision}
\end{align}
Let $v_{0, P}[u, l]$ denote the mixture prior variance. The Gaussian mixture test supermartingale for precision is:
\begin{equation}\label{eq:gm_precision}
M_t^{(P)}(u, l) \coloneqq \sqrt{\frac{v_{0, P}[u, l]}{v_{0, P}[u, l] + V_t^{(P)}(u, l)}} \exp\left( \frac{\big(\max\{0, S_t^{(P)}(u, l)\}\big)^2}{2\big(v_{0, P}[u, l] + V_t^{(P)}(u, l)\big)} \right).
\end{equation}
To control the family-wise error rate across all $K_l$ candidate lower thresholds simultaneously, precision tests apply a Bonferroni correction over lower thresholds, rejecting $H_0^{(P)}(u, l)$ when:
\begin{equation}\label{eq:reject_precision}
M_t^{(P)}(u, l) \ge \frac{1}{\alpha_P} \quad \text{with } \alpha_P \coloneqq \frac{\delta_P}{K_l} = \frac{\delta}{2 K_l}.
\end{equation}

\subsection{Threshold Selection and the Active Martingale Pair}
At any stopping time $\tau_{\text{stop}}$, ACTIS determines the active certified threshold pair $(l^*, u^*)$ via:
\begin{align}
l^*(\tau_{\text{stop}}) &\coloneqq \max \Big\{ l \in \{0, \dots, K_l-1\} : M_{\tau_{\text{stop}}}^{(R)}(l) \ge \frac{1}{\delta_R} \Big\}, \label{eq:select_l} \\
u^*(\tau_{\text{stop}}) &\coloneqq \min \Big\{ u \in \{l^*, \dots, K_u-1\} : M_{\tau_{\text{stop}}}^{(P)}(u, l^*) \ge \frac{1}{\alpha_P} \Big\}. \label{eq:select_u}
\end{align}
If no lower threshold passes the recall test, the system defaults to the deterministic fallback anchor $l^* = 0$ ($\tau_0 = 0.0$), which guarantees $100\%$ recall by construction. If no upper threshold passes the precision test at $l^*$, the system defaults to the deterministic fallback anchor $u^* = K_u - 1$ ($\tau_{K_u-1}' = 1.0$), which guarantees $100\%$ precision by construction.

\begin{remark}[\textbf{Specification of the Evaluated Martingales}]\label{rem:which_martingale}
A critical ambiguity in earlier diagnostic discussions was the failure to specify which martingales were being tested. In a grid with $K_l = 20$ and $K_u = 20$, there exist $20$ recall martingales $\{M_t^{(R)}(l)\}_{l=0}^{19}$ and up to $210$ precision martingales $\{M_t^{(P)}(u, l)\}_{0 \le l \le u < 20}$. 

Demanding that every latent martingale simultaneously satisfy CLT variance accumulation is mathematically unnecessary and statistically flawed: the majority of thresholds are never accepted. The statistical guarantee of the cascade~\eqref{eq:joint_guarantee} depends strictly on the validity of the \textbf{active threshold pair} $(l^*, u^*)$ selected at stopping. Thus, asymptotic validity diagnostics must be evaluated specifically on:
\begin{equation}
\big(V_t^{(R)}(l^*), \; v_{0, R}[l^*]\big) \quad \text{and} \quad \big(V_t^{(P)}(u^*, l^*), \; v_{0, P}[u^*, l^*]\big).
\end{equation}
If $l^* = 0$ or $u^* = K_u - 1$, the threshold operates at a deterministic fallback anchor, which is error-free by construction and completely exempt from statistical normal approximations.
\end{remark}

\newpage
\section{Foundations of the Martingale Central Limit Theorem}

The validity of Ville's inequality~\eqref{eq:ville_ineq} for the Gaussian mixture wealth process~\eqref{eq:gm_recall} relies on the assumption that the log-likelihood ratio increments are sub-Gaussian. When increment variables $X_t$ are bounded non-Gaussian random variables, Robbins~\cite{robbins1970} and Howard et al.~\cite{howard2021} showed that the mixture boundary crossing probabilities hold asymptotically as the accumulated variance grows large. 

However, because the calibration procedure terminates at a data-dependent stopping time $\tau_{\text{stop}}$, the classical independent and identically distributed (i.i.d.) Central Limit Theorem is inapplicable. The sequential stopping process is governed strictly by the \textbf{Martingale Central Limit Theorem}.

\subsection{The Martingale Limit Setting}
Let $(\Omega, \mathcal{F}, (\mathcal{F}_t)_{t \ge 0}, \mathbb{P})$ be a filtered probability space. Let $\{D_t, \mathcal{F}_t\}_{t \ge 1}$ be a martingale difference sequence such that $\mathbb{E}[D_t \mid \mathcal{F}_{t-1}] = 0$ and $\mathbb{E}[D_t^2] < \infty$. The partial sum process $S_t = \sum_{i=1}^t D_i$ is a zero-mean martingale.

We define two fundamental variance processes:
\begin{align}
\text{Predictable quadratic variation: } \langle S \rangle_t &\coloneqq \sum_{i=1}^t \mathbb{E}\big[ D_i^2 \;\big|\; \mathcal{F}_{i-1} \big], \label{eq:angle_bracket} \\
\text{Realized quadratic variation: } [S]_t &\coloneqq \sum_{i=1}^t D_i^2. \label{eq:square_bracket}
\end{align}
In the centered sample variance accumulator $V_t = \sum_{i=1}^t (D_i - \bar{D}_t)^2$, the identity $V_t = [S]_t - t \bar{D}_t^2$ holds. Under the null hypothesis where $\bar{D}_t \xrightarrow{p} 0$, $V_t / [S]_t \xrightarrow{p} 1$ as $t \to \infty$.

\subsection{Authoritative Martingale CLT Theorems}
The foundational theory of central limit theorems for martingales under random norming and stopping times was established by Brown (1971)~\cite{brown1971}, McLeish (1974)~\cite{mcleish1974}, and consolidated in the authoritative treatise by \textbf{Hall and Heyde (1980)}~\cite{hall1980martingale}.

\begin{theorem}[\textbf{Martingale Central Limit Theorem under Stopping Times; Hall \& Heyde 1980, Theorem 3.4}]\label{thm:mclt}
Let $\{S_{t}, \mathcal{F}_{t}\}_{t \ge 1}$ be a zero-mean martingale with differences $D_t$, and let $\{\tau_m\}_{m \ge 1}$ be a sequence of stopping times such that $\tau_m \to \infty$ almost surely as $m \to \infty$. Let $s_m^2$ be a sequence of positive scaling constants (or $\mathcal{F}_{\tau_m}$-measurable random variables) such that $s_m^2 \to \infty$. 

Suppose the following two conditions are satisfied:
\begin{enumerate}[label=(\alph*)]
    \item \textbf{Quadratic Variation Consistency}:
    \begin{equation}\label{eq:cond_quad_var}
    \frac{[S]_{\tau_m}}{s_m^2} \xrightarrow{p} \eta^2 \quad (\text{or } \frac{\langle S \rangle_{\tau_m}}{s_m^2} \xrightarrow{p} \eta^2),
    \end{equation}
    where $\eta^2$ is an almost surely positive, finite random variable.
    \item \textbf{Conditional Lindeberg Condition}: For every $\epsilon > 0$,
    \begin{equation}\label{eq:cond_lindeberg}
    \frac{1}{s_m^2} \sum_{i=1}^{\tau_m} \mathbb{E}\Big[ D_i^2 \mathbf{1}_{\{|D_i| > \epsilon s_m\}} \;\Big|\; \mathcal{F}_{i-1} \Big] \xrightarrow{p} 0 \quad \text{as } m \to \infty.
    \end{equation}
\end{enumerate}
Then, the normalized stopped martingale converges in distribution to a variance mixture of normals:
\begin{equation}
\frac{S_{\tau_m}}{s_m} \xrightarrow{d} Z \cdot \eta,
\end{equation}
where $Z \sim \mathcal{N}(0, 1)$ is independent of $\eta$. In particular, if $\eta^2 = 1$ in probability, then:
\begin{equation}
\frac{S_{\tau_m}}{\sqrt{[S]_{\tau_m}}} \xrightarrow{d} \mathcal{N}(0, 1).
\end{equation}
\end{theorem}

\subsection{Implications for Anytime Boundary Crossing}
In sequential analysis, we are not merely interested in the distribution of $S_t$ at a single fixed $t$, but in the probability of crossing a continuous curve over an open horizon: $\mathbb{P}(\exists t \ge 1 : S_t \ge h(V_t))$. 

Lai (1976)~\cite{lai1976boundary} and Robbins (1970)~\cite{robbins1970} established that boundary crossing probabilities for discrete-time martingales converge to the boundary crossing probabilities of standard Brownian motion if the discrete increments $D_t$ are uniformly negligible relative to the total accumulated variance. 

Specifically, this requires the \textbf{Maximum Jump-to-Variance Condition}:
\begin{equation}\label{eq:max_jump_ratio}
\rho_{\max}(t) \coloneqq \frac{\max_{1 \le i \le t} D_i^2}{[S]_t} \xrightarrow{p} 0 \quad \text{as } [S]_t \to \infty.
\end{equation}

\begin{remark}[\textbf{Origin of the Diagnostic Desiderata}]\label{rem:desiderata_origin}
The user rightly asked: \emph{``Where have these desiderata come from? Shouldn't we be citing some authoritative source?''}

The three runtime criteria implemented in our refined ACTIS diagnostic are derived directly from the foundational conditions of Theorem~\ref{thm:mclt} and Lai's boundary crossing invariance principle:
\begin{enumerate}[leftmargin=*]
    \item \textbf{Condition (a) (Quadratic Variation Accumulation)}: Requires that realized quadratic variation $V_t$ has reached a non-trivial fraction of the expected process scale ($V_t \ge c \cdot v_{0, \text{target}}$), ensuring that the process is operating in a regime where variance has stabilized.
    \item \textbf{Condition (b) (Lindeberg Dispersion / Jump Bound)}: Measured via the jump ratio $\rho_{\max}(t) \le \kappa$. If a single observation contributes $60\%$ or $90\%$ of the total empirical variance, the conditional Lindeberg condition is severely violated, and the tail distribution remains discrete and non-Gaussian.
    \item \textbf{Berry-Esseen Skewness Floor}: When observations are driven by sparse binary events ($Y_i = 1$), the third-moment error of the normal approximation is bounded by $\frac{C \mathbb{E}[|D_i|^3]}{\sigma^3 \sqrt{t}} \sim \frac{1}{\sqrt{n_{\text{pos}}}}$. Requiring $n_{\text{pos}} \ge \text{min\_positives} \approx 8 \log(2/\delta)$ guarantees that higher-order skewness has dissipated.
\end{enumerate}
\end{remark}

\newpage
\section{Critical Appraisal of Kish's Effective Sample Size}

In initial diagnostic explorations, it was suggested that importance weights could be accounted for by applying Kish's Effective Sample Size (ESS) formula to scale the effective number of observations. The user asked:
\begin{center}
\emph{``Is Kish's effective sample size typically applied to martingales? And why are importance weights dropped for the diagnostics?''}
\end{center}

We present a critical, first-principles examination of this question.

\subsection{Derivation and Domain of Kish's Formula}
In his classical 1965 text \emph{Survey Sampling}~\cite{kish1965survey}, Leslie Kish investigated the variance inflation that occurs when estimating a population mean $\bar{Y} = \frac{1}{N}\sum_{i=1}^N Y_i$ from a non-uniformly weighted sample. Let $w_1, \dots, w_n$ denote fixed, normalized survey weights assigned to $n$ independent observations. The weighted estimator is $\bar{y}_w = \sum_{i=1}^n w_i Y_i / \sum_{i=1}^n w_i$.

Kish defined the \textbf{Design Effect} ($\text{deff}$) as the ratio of the variance of the weighted estimator to that of an unweighted Simple Random Sample (SRS) of equal size:
\begin{equation}\label{eq:kish_deff}
\text{deff} \coloneqq \frac{\text{Var}(\bar{y}_w)}{\text{Var}_{\text{SRS}}(\bar{y})} \approx 1 + \text{cv}_w^2 = 1 + \frac{\text{Var}(w)}{\bar{w}^2} = \frac{n \sum_{i=1}^n w_i^2}{(\sum_{i=1}^n w_i)^2}.
\end{equation}
The \textbf{Kish Effective Sample Size} ($n_{\text{eff}}$) represents the number of unweighted SRS observations that would yield equivalent variance:
\begin{equation}\label{eq:kish_ess}
n_{\text{eff}} \coloneqq \frac{n}{\text{deff}} = \frac{\big(\sum_{i=1}^n w_i\big)^2}{\sum_{i=1}^n w_i^2}.
\end{equation}

\subsection{Why Kish's ESS Does Not Apply to Test Martingales}
Kish's derivation is fundamentally rooted in \textbf{static survey sampling of finite populations} under fixed sample sizes. It is \textbf{not} typically applied to, nor mathematically formulated for, sequential test martingales. There are three profound mathematical reasons why applying Kish's ESS to martingale diagnostics is inappropriate:

\begin{enumerate}[leftmargin=*]
    \item \textbf{Quadratic Variation Already Contains the Squared Weights}:
    In martingale theory, the variance of the partial sum process $S_t = \sum_{i=1}^t X_i$ with $X_i = w(Z_i) h(Z_i, Y_i)$ is tracked \emph{directly and exactly} by the sample variance accumulator:
    \begin{equation}\label{eq:vt_exact}
    V_t = \sum_{i=1}^t \big( X_i - \bar{X}_t \big)^2 = \sum_{i=1}^t \Big( w(Z_i) h(Z_i, Y_i) - \bar{X}_t \Big)^2.
    \end{equation}
    Notice that the importance weights $w(Z_i)$ appear \textbf{inside the quadratic sum as $w(Z_i)^2$}. 
    
    The actual dispersion and variance inflation caused by unequal importance weights is already explicitly accumulated in $V_t$. If one computes $V_t$ and then additionally scales the sample size or variance by Kish's $\text{deff}$, one is penalizing the weights \textbf{twice} ($w_i^4$ scaling). In martingale limit theory, Theorem~\ref{thm:mclt} requires convergence of $V_t$ itself; there is no role for a surrogate SRS equivalent.
    
    \item \textbf{Optimal Importance Proposals Reduce Variance, Violating Kish's Assumption}:
    Kish's formula assumes that unequal weighting is a statistical nuisance (e.g. post-stratification corrections) that strictly inflates variance relative to uniform sampling ($\text{deff} \ge 1$). 
    
    In contrast, in importance sampling for rare events, the proposal distribution $q(x) \propto s(x)$ is chosen specifically to \emph{reduce} the variance of the indicator $Y_i \mathbf{1}\{s(x_i) \ge \tau\}$ by oversampling the rare region where $Y_i = 1$. When the proposal is well-aligned with the oracle signal, the actual variance $\text{Var}_q(w(Z) Y)$ is substantially \emph{smaller} than the unweighted variance $\text{Var}_{\text{unif}}(Y)$. Applying Kish's formula penalizes the efficient proposal for having unequal weights, misidentifying a variance-reducing sampling scheme as a low-information sample.
    
    \item \textbf{Distinction from Sequential Monte Carlo (SMC)}:
    In particle filtering and SMC algorithms (e.g., Doucet et al.~\cite{doucet2000sequential}), Kish's ESS is frequently computed across particles: $\text{ESS}_t = 1 / \sum_{j=1}^P (W_t^{(j)})^2$. However, in SMC, $\text{ESS}_t$ is evaluated across parallel particles at a fixed time step to detect particle degeneracy and trigger resampling. It is not used as a boundary-crossing test martingale, nor does it replace the quadratic variation of a stopped temporal process.
\end{enumerate}

\subsection{Why Weights are Separated from Event Counts in Diagnostics}
This clarifies why our diagnostics do not apply importance weights to the raw positive count $n_{\text{pos}} = \sum_{i=1}^t Y_i$:
\begin{itemize}[leftmargin=*]
    \item The role of $V_t$ is to measure the \textbf{process variance scale}, and it fully incorporates the weights $w(Z_i)^2$.
    \item The role of $n_{\text{pos}}$ is to bound the \textbf{higher-moment skewness} of the discrete Bernoulli observations $Y_i \in \{0, 1\}$ via the Berry-Esseen theorem. Bounding the third centralized moment $\mathbb{E}[|Y_i - p|^3]$ depends on observing a minimum physical number of true positive labels from the oracle. Weighting $n_{\text{pos}}$ by $w_i$ would distort the count of physical Bernoulli trials.
\end{itemize}

\newpage
\section{Deconstruction of the Diagnostic Failure in ACTIS}

We now analyze the exact mathematical bug in \texttt{AsymptoticValidityDiagnostics} that led to the severe performance degradation observed by the user.

\subsection{The Original Implementation}
In the existing ACTIS codebase, the diagnostic method evaluated the ratio of accumulated empirical variance to the effective mixture prior variance:
\begin{equation}\label{eq:orig_ratio}
\text{ratio\_recall} \coloneqq \frac{V_t^{(R)}(l^*)}{v_{0, \text{recall}}[l^*]} \ge 0.05.
\end{equation}
If $\text{ratio\_recall} < 0.05$ (or $\text{ratio\_precision} < 0.05$), the diagnostic returned $\texttt{asymp\_valid} = \text{False}$. The tuner interpreted this as indicating that pre-asymptotic warmup was incomplete, setting $\texttt{warmup\_needed} = \text{True}$ and forcing the sampling loop to continue drawing queries.

\subsection{The Zero-Variance Floor and Asymmetric Jump Bounds}
Recall the definition of the threshold-localized prior variance from Section 4 of the technical specification:
\begin{equation}\label{eq:v0_recall_def}
v_{0, \text{recall}}[l] \coloneqq \max\Big( v_{0, \text{floor}}^{(R)}[l], \; n_0 \cdot \hat{\sigma}_R^2[l] \Big).
\end{equation}
The zero-variance safety floor was derived in Section 3 of the technical specification:
\begin{equation}\label{eq:v0_floor_formula}
v_{0, \text{floor}}^{(R)}[l] = \frac{2 b_R^2[l]}{\log(2/\delta)},
\end{equation}
where $b_R[l]$ bounds the maximum possible increment in absolute value:
\begin{equation}\label{eq:bR_formula}
b_R[l] = \max\Big( (1 - \gamma_R) w_{\max, \ge \tau_l}, \; \gamma_R w_{\max, < \tau_l} \Big).
\end{equation}

Under high target recall constraints (such as $\gamma_R = 0.95$), notice the massive asymmetry between positive and negative jumps:
\begin{align}
\text{Positive jump (True Positive):} \quad & X_{t, l}^{(R)} = (1 - \gamma_R) w(Z_t) = 0.05 \cdot w(Z_t), \\
\text{Negative jump (False Negative):} \quad & X_{t, l}^{(R)} = -\gamma_R w(Z_t) = -0.95 \cdot w(Z_t).
\end{align}
The negative jump is nineteen times larger than the positive jump! Consequently, the maximum jump bound $b_R[l]$ is completely dominated by the negative jump:
\begin{equation}
b_R[l] = \gamma_R w_{\max, < \tau_l} \approx 0.95 \times 2.40 = 2.28.
\end{equation}
For $\delta = 0.20$, $\log(2 / \delta) = \log(10) \approx 2.3026$. The resulting safety floor is:
\begin{equation}
v_{0, \text{floor}}^{(R)}[l] = \frac{2 \times (2.28)^2}{2.3026} \approx 4.515.
\end{equation}

\subsection{The Mechanism of Failure on High-Recall Paths}
Now consider what occurs when ACTIS evaluates an accurate candidate lower threshold $\tau_l^*$. Because $\tau_l^*$ satisfies the high recall constraint ($\text{Recall}(\tau_l^*) \ge 0.95$), in any typical calibration sample of $t = 500\text{--}1000$ items, \textbf{the observed number of false negatives below $\tau_l^*$ is zero} ($n_{\text{FN}} = 0$).

Every positive item sampled by the tuner falls above $\tau_l^*$. Therefore, every non-zero increment is strictly positive:
\begin{equation}
X_{i, l^*}^{(R)} = (1 - \gamma_R) w(Z_i) = 0.05 \cdot w(Z_i).
\end{equation}
Items with negative labels ($Y_i = 0$) contribute identically zero: $X_{i, l^*}^{(R)} = 0$.

Let $n_{\text{pos}} = \sum_{i=1}^t Y_i$ denote the number of positive items observed in the calibration sample. The sample variance accumulator is:
\begin{equation}\label{eq:vt_scaling}
V_t^{(R)}(l^*) = \sum_{i=1}^t \big(X_{i, l^*}^{(R)} - \bar{X}_t^{(R)}\big)^2 \approx n_{\text{pos}} \cdot (1 - \gamma_R)^2 \cdot \bar{w}_{\text{pos}}^2.
\end{equation}
Crucially, $V_t^{(R)}(l^*)$ scales with \textbf{$(1 - \gamma_R)^2 = (0.05)^2 = 0.0025$}.

Suppose the tuner has drawn an initial sample containing $n_{\text{pos}} = 30$ true positives with typical weights $w \approx 1.2$. The accumulated sample variance is:
\begin{equation}
V_t^{(R)}(l^*) \approx 30 \times (0.05 \times 1.2)^2 = 30 \times 0.0036 = 0.108.
\end{equation}
Now, compare this accumulated variance to the safety floor $v_{0, \text{floor}}^{(R)} \approx 4.515$:
\begin{equation}
\frac{V_t^{(R)}(l^*)}{v_{0, \text{recall}}[l^*]} = \frac{0.108}{4.515} \approx 0.0239 < 0.05.
\end{equation}

\begin{theorem}[\textbf{Dimensional Scale Mismatch}]\label{thm:mismatch}
Let $\gamma_R \in (0, 1)$ be the target recall and suppose no false negatives have been observed in the sample ($n_{\text{FN}} = 0$). Then the ratio of the empirical variance accumulator $V_t^{(R)}$ to the zero-variance safety floor $v_{0, \text{floor}}^{(R)}$ scales as:
\begin{equation}
\frac{V_t^{(R)}}{v_{0, \text{floor}}^{(R)}} \propto \left( \frac{1 - \gamma_R}{\gamma_R} \right)^2.
\end{equation}
For $\gamma_R = 0.95$, $((1 - \gamma_R)/\gamma_R)^2 = (0.05 / 0.95)^2 = 1 / 361 \approx 0.00277$.
\end{theorem}

To satisfy $V_t^{(R)} / v_{0, \text{floor}} \ge 0.05$, the tuner would require:
\begin{equation}
V_t^{(R)} \ge 0.05 \times 4.515 = 0.2258 \implies n_{\text{pos}} \ge \frac{0.2258}{0.0036} \approx 63 \text{ positive items}.
\end{equation}

On real-world datasets, this requirement is disastrous:
\begin{itemize}[leftmargin=*]
    \item On \textbf{SUPG OntoNotes}, there are only 279 total positive items in the entire database of 11,165 records. Capturing 63 positives requires sampling more than $2,500$ records in calibration alone.
    \item On \textbf{SUPG ImageNet}, there are only 50 positive items in the entire dataset of 50,000 records. It is physically impossible to observe 63 positives even if the entire database is exhausted!
\end{itemize}

As a result, \texttt{ratio\_recall} $< 0.05$ held permanently throughout the entire run. The diagnostic systematically declared \texttt{asymp\_valid = False}, overriding the value-of-information stopping rule and forcing ACTIS to sample continuously until hitting $\text{max\_sample\_size} = 0.50 N$.

\subsection{Why Ad-Hoc Branching on $n_{\text{FN}} = 0$ is Unprincipled}
In preliminary discussions, an ad-hoc fix was considered:
\begin{quote}
\emph{``If $n_{\text{FN}} == 0$, exempt the variance ratio and check only $n_{\text{pos}} \ge 10$.''}
\end{quote}
The user rightly expressed skepticism:
\begin{quote}
\emph{``Is it sound for the behavior to change based on the observed number of FN and FP? Or is there a theoretical reason why the current approach is flawed?''}
\end{quote}

We strongly agree with the user's critique. Branching conditionally on $n_{\text{FN}} = 0$ is theoretically unprincipled for two reasons:
\begin{enumerate}[leftmargin=*]
    \item \textbf{Post-Selection Distortion}: $n_{\text{FN}}$ is a random, endogenous quantity of the sample path. Creating a discontinuous bifurcation in the diagnostic boundary based on whether a rare event occurred distorts the stopping distribution and lacks any justification under martingale limit theory.
    \item \textbf{Misidentification of the Cause}: The failure was \textbf{not} that CLT convergence behaves discontinuously when $n_{\text{FN}} = 0$. The failure was that the original diagnostic compared $V_t$ to $v_{0, \text{floor}}$ instead of $v_{0, \text{target}}$! 
\end{enumerate}
$v_{0, \text{floor}}$ was derived in Section 3 of the technical specification strictly for Theorem 1: to guarantee mathematically that an uninformative run of $t \le 2$ observations cannot reject the null under zero empirical variance. It is a non-asymptotic safety clamp against the worst-case negative jump for $t \le 2$. It was \textbf{never intended} to serve as an indicator of the empirical process variance scale at $t \gg 2$.

\newpage
\section{The Principled Diagnostic Framework}

Having identified the root theoretical cause, we formulate a principled, unified diagnostic framework derived directly from Martingale Central Limit Theorem conditions.

\subsection{Expected Cumulative Process Variance $v_{0, \text{target}}$}
To evaluate whether empirical variance has accumulated sufficiently, $V_t$ must be compared to the \textbf{expected cumulative process variance} over the calibration horizon $n_0$:
\begin{align}
v_{0, \text{target}}^{(R)}[l^*] &\coloneqq n_0 \cdot \hat{\sigma}_R^2[l^*], \label{eq:v0_target_recall} \\
v_{0, \text{target}}^{(P)}[u^*, l^*] &\coloneqq n_0 \cdot \hat{\sigma}_P^2[u^*, l^*], \label{eq:v0_target_precision}
\end{align}
where $\hat{\sigma}_R^2[l^*]$ and $\hat{\sigma}_P^2[u^*, l^*]$ are the population proxy surrogate variances defined in Equations~(164)--(165) of the technical specification:
\begin{align}
\hat{\sigma}_R^2[l] &\coloneqq \frac{(1 - \gamma_R)^2}{N}\sum_{i: s_i \ge \tau_l} s_i w_i + \frac{\gamma_R^2}{N}\sum_{i: s_i < \tau_l} s_i w_i, \\
\hat{\sigma}_P^2[u, l] &\coloneqq \frac{\gamma_P^2}{N}\sum_{i: s_i \ge \tau_u'} (1 - s_i) w_i + \frac{(1 - \gamma_P)^2}{N}\sum_{i: s_i \ge \tau_l} s_i w_i.
\end{align}

Notice that when $\tau_l^*$ is an accurate threshold (where false negatives $\sum_{s < \tau_l} s_i w_i \approx 0$), $\hat{\sigma}_R^2[l^*]$ naturally operates on the $(1 - \gamma_R)^2$ scale! Both $V_t^{(R)}(l^*)$ and $v_{0, \text{target}}^{(R)}[l^*]$ share the exact same physical scale and dimension.

\subsection{Unified Diagnostic Specification}
Let $(l^*, u^*)$ denote the candidate threshold pair proposed for certification at stopping step $t$. The thresholds are certified as asymptotically valid if and only if:
\begin{enumerate}[leftmargin=*]
    \item \textbf{Deterministic Fallback Anchor Exemption}:
    If $l^* = 0$ ($\tau_{\text{neg}} = 0.0$), the recall test is deterministically error-free ($\text{Recall} \equiv 1.0$), exempting recall variance checks. If $u^* = K_u - 1$ ($\tau_{\text{pos}} = 1.0$), the precision test is deterministically error-free ($\text{Precision} \equiv 1.0$), exempting precision variance checks.
    
    \item \textbf{Quadratic Variation Accumulation (Condition~\ref{eq:cond_quad_var})}:
    For non-anchor thresholds, empirical variance must reach at least $5\%$ of the expected process variance:
    \begin{align}
    V_t^{(R)}(l^*) &\ge 0.05 \cdot v_{0, \text{target}}^{(R)}[l^*], \label{eq:check_var_r} \\
    V_t^{(P)} (u^*, l^*) &\ge 0.05 \cdot v_{0, \text{target}}^{(P)}[u^*, l^*]. \label{eq:check_var_p}
    \end{align}
    
    \item \textbf{Hall--Heyde Conditional Lindeberg Jump Bound (Condition~\ref{eq:cond_lindeberg})}:
    To ensure that no single jump dominates the accumulated variance, the maximum jump-to-variance ratio must satisfy:
    \begin{align}
    \rho_{\max}^{(R)}(t) &\coloneqq \frac{\max_{1 \le i \le t} \big(X_{i, l^*}^{(R)} - \bar{X}_{t, l^*}^{(R)}\big)^2}{V_t^{(R)}(l^*)} \le \kappa, \label{eq:check_jump_r} \\
    \rho_{\max}^{(P)}(t) &\coloneqq \frac{\max_{1 \le i \le t} \big(X_{i, u^*, l^*}^{(P)} - \bar{X}_{t, u^*, l^*}^{(P)}\big)^2}{V_t^{(P)}(u^*, l^*)} \le \kappa, \label{eq:check_jump_p}
    \end{align}
    where $\kappa = 0.50$ guarantees that no single observation accounts for more than half of the accumulated variance.
    
    \item \textbf{Berry-Esseen Non-Degenerate Event Support}:
    Total observed positive items $n_{\text{pos}} \ge \text{min\_positives} \coloneqq \min(N, \min(\lceil 8 \log(2/\delta) \rceil, \max(5, \lceil p_{\text{eff}} N \rceil)))$, and true positive support above $\tau_{l^*}$ satisfies $n_{\text{TP}} \ge 5$.
\end{enumerate}

This formulation contains \textbf{zero ad-hoc branching on $n_{\text{FN}}$ or $n_{\text{FP}}$}. It applies uniformly across all sample paths, candidate thresholds, and data distributions.

\subsection{Necessity and Non-Redundancy of the Diagnostic Conditions}
The four conditions form a logical conjunction ($C_1 \land C_2 \land C_3 \land C_4$); all applicable checks must be satisfied simultaneously. They are \textbf{not mutually exclusive}, nor are they redundant. Each check addresses a fundamentally distinct failure mode of asymptotic inference:

\begin{enumerate}[leftmargin=*]
    \item \textbf{Condition 1 (Anchor Exemption)}: Applies strictly to deterministic boundary anchors ($l^* = 0$ or $u^* = K_u - 1$). Because $\text{Recall}(0.0) \equiv 1.0$ and $\text{Precision}(1.0, \tau_{\text{neg}}) \equiv 1.0$ hold deterministically without estimation, they are exempt from normal approximation requirements.
    \item \textbf{Condition 2 (Quadratic Variation Scale $V_t \ge 0.05 v_{0, \text{target}}$)}: Protects against the \emph{pre-asymptotic zero-variance regime}. Hall and Heyde (1980, Theorem 3.4) require that predictable quadratic variation does not vanish relative to nominal variance ($[S]_t / s_t^2 \xrightarrow{p} \eta^2 > 0$). If $V_t \approx 0$, the normal approximation is mathematically undefined.
    \begin{itemize}
        \item \emph{Why Condition 3 cannot replace Condition 2}: Suppose $t = 30$ samples have tiny, nearly identical values, yielding $V_t = 10^{-12}$. The jump dispersion ratio $\rho_{\max} = 1/30 = 0.033 \le 0.50$ easily passes, but the process has accumulated zero usable variance. Condition 2 detects this deficiency.
    \end{itemize}
    \item \textbf{Condition 3 (Hall--Heyde Conditional Lindeberg Jump Bound $\rho_{\max} \le 0.50$)}: Protects against \emph{heavy-tailed jump dominance / single-point variance concentration}. Hall and Heyde (1980, Theorem 3.2 and Section 3.2) explicitly state that the conditional Lindeberg condition requires individual terms to be uniformly negligible relative to total variance.
    \begin{itemize}
        \item \emph{Why Condition 2 cannot replace Condition 3}: Suppose $V_t = 10.0 \ge 0.05 v_{0, \text{target}}$ is large. But $V_t$ consists of $99$ observations with increment $0$ and \textbf{one single observation with a large importance weight} $w_i = 60$ contributing $90\%$ of the variance ($\rho_{\max} = 0.90$). The process is a single-jump step function, not a continuous Gaussian path! Condition 3 catches this failure mode.
    \end{itemize}
    \item \textbf{Condition 4 (Berry-Esseen Event Support)}: Protects against \emph{higher-moment skewness in small-sample discrete Bernoulli trials}.
    \begin{itemize}
        \item \emph{Why Conditions 2 and 3 cannot replace Condition 4}: The Martingale CLT is strictly an \emph{asymptotic} theorem ($t \to \infty$); it does not provide non-asymptotic uniform rates of convergence. For rare binary events ($Y_i \in \{0, 1\}$ with prevalence $p \ll 1$), Berry-Esseen and Edgeworth expansions show that approximation error scales as $O(1/\sqrt{n_{\text{pos}}})$. If $n_{\text{pos}} = 2$, $\rho_{\max}$ can mathematically be $0.50$ and $V_t$ can exceed a small $v_{0, \text{target}}$, but a sum of two Bernoulli trials has an overwhelmingly discrete, skewed distribution where Gaussian tail probabilities deviate substantially from true coverage.
    \end{itemize}
\end{enumerate}

\subsection{Mathematical Connection of Event Counts to the Active Martingales}
The event count thresholds in Condition 4 ($n_{\text{pos}}$ and $n_{\text{TP}}$) are not arbitrary global counts; they connect directly to the mathematical definitions of the active martingales at candidate threshold pair $(l^*, u^*)$:

\subsubsection{Connection to the Active Recall Martingale $M_t^{(R)}(l^*)$}
The sequential increment is:
\begin{equation}
X_{t, l^*}^{(R)} = w(Z_t) Y_t \big( \mathbf{1}\{s(Z_t) \ge \tau_{l^*}\} - \gamma_R \big).
\end{equation}
Whenever ground-truth label $Y_t = 0$, $X_{t, l^*}^{(R)} = 0$ identically. Ground-truth negative items provide \textbf{zero signal and zero uncentered variance} to the recall martingale. Non-zero increments occur if and only if $Y_t = 1$:
\begin{itemize}
    \item \textbf{True Positives ($s \ge \tau_{l^*}$)}: Count $n_{\text{TP}}(l^*)$, each contributing positive increment $(1 - \gamma_R) w(Z_t) > 0$.
    \item \textbf{False Negatives ($s < \tau_{l^*}$)}: Count $n_{\text{FN}}(l^*)$, each contributing negative increment $-\gamma_R w(Z_t) < 0$.
\end{itemize}
Therefore, the total number of informative observations in the recall martingale is:
\begin{equation}
n_{\text{info}}^{(R)}(l^*) = n_{\text{TP}}(l^*) + n_{\text{FN}}(l^*) = \sum_{i=1}^t Y_i = n_{\text{pos}}.
\end{equation}
Thus, $n_{\text{pos}}$ is the \textbf{exact number of active increments in the recall martingale difference sequence}. Bounding Berry-Esseen skewness below $\delta/2$ requires $n_{\text{pos}} \ge 8 \log(2/\delta)$.

\subsubsection{Connection to the Active Precision Martingale $M_t^{(P)}(u^*, l^*)$}
The sequential increment is:
\begin{equation}
X_{t, u^*, l^*}^{(P)} = w(Z_t) \Big( (1 - \gamma_P) Y_t \mathbf{1}\{s(Z_t) \ge \tau_{l^*}\} - \gamma_P (1 - Y_t) \mathbf{1}\{s(Z_t) \ge \tau_{u^*}\} \Big).
\end{equation}
Non-zero increments come from two distinct events:
\begin{enumerate}
    \item True Positives routed to oracle or cascade output: $Y_t = 1$ and $s(Z_t) \ge \tau_{l^*}$ (count $n_{\text{TP}}(l^*)$), contributing $(1 - \gamma_P) w(Z_t) > 0$.
    \item False Positives automated without oracle review: $Y_t = 0$ and $s(Z_t) \ge \tau_{u^*}$ (count $n_{\text{FP}}(u^*)$), contributing $-\gamma_P w(Z_t) < 0$.
\end{enumerate}
Crucially, \textbf{$n_{\text{TP}}(l^*)$ provides the entire positive wealth growth} for the precision martingale! If $n_{\text{TP}}(l^*) < 5$, the cascade has observed almost zero positive classifications. A cascade cannot reliably certify that its precision is $\ge \gamma_P$ if it has only encountered 1 or 2 positive predictions. Requiring $n_{\text{TP}}(l^*) \ge 5$ ensures empirical positive evidence for the cascade's routing band.

\newpage
\section{Empirical Validation and Status Across the SUPG Benchmarks}

\subsection{Experimental Scope and Status}
To assess the impact of the refined asymptotic diagnostics, we study performance across the four benchmark datasets in the SUPG suite~\cite{kang2020approximate}:
\begin{enumerate}[leftmargin=*]
    \item \textbf{SUPG OntoNotes} ($N = 11,165$, prevalence $2.50\%$): Named entity recognition filtering.
    \item \textbf{SUPG TACRED} ($N = 22,631$, prevalence $2.36\%$): Relation extraction filtering.
    \item \textbf{SUPG Jackson} ($N = 973,085$, prevalence $29.24\%$): Large-scale video object detection.
    \item \textbf{SUPG ImageNet} ($N = 50,000$, prevalence $0.10\%$): Extreme rare-event image classification ($50$ positives in $50,000$).
\end{enumerate}

\paragraph{Empirical Status Note:} The results presented below reflect full Slurm cluster array execution (Job ID \texttt{30259309}, 20 paired Monte Carlo trials per configuration) across all four benchmarks in the SUPG suite under the unified four-condition diagnostic implemented in the codebase.

\subsection{Benchmark Performance at $\delta = 0.20$}
Table~\ref{tab:results_delta_20} reports empirical performance under $\delta = 0.20$.

\begin{table}[ht]
\centering
\small
\setlength{\tabcolsep}{5pt}
\caption{Empirical benchmark comparison at $\gamma_R = \gamma_P = 0.95$, $\delta = 0.20$ (20 trials).}\label{tab:results_delta_20}
\begin{tabular*}{\textwidth}{@{\extracolsep{\fill}}llrrrr}
\toprule
Dataset & Method & Oracle Calls & Fail Rate & Mean $n_{\text{calib}}$ & Clean Stop \\
\midrule
OntoNotes & BARGAIN-PR & 42.93\% & 5.0\% & 1,675 & 100\% \\
($N=11,165$) & ACTIS (WOR) & 38.88\% & 0.0\% & 1,116 & 75\% \\
 & ACTIS IS (Original) & 47.56\% & 0.0\% & 3,564 & 0\% \\
 & \textbf{ACTIS IS (Principled)} & \textbf{42.26\%} & \textbf{0.0\%} & \textbf{2,830} & \textbf{100\%} \\
\midrule
TACRED & BARGAIN-PR & 40.93\% & \textbf{10.0\%} & 2,263 & 100\% \\
($N=22,631$) & ACTIS (WOR) & 35.53\% & \textbf{10.0\%} & 2,263 & 100\% \\
 & ACTIS IS (Original) & 49.32\% & 0.0\% & 7,242 & 0\% \\
 & \textbf{ACTIS IS (Principled)} & \textbf{36.57\%} & \textbf{0.0\%} & \textbf{3,421} & \textbf{95\%} \\
\midrule
Jackson & BARGAIN-PR & 76.87\% & \textbf{15.0\%} & 3,000 & 100\% \\
($N=973,085$) & ACTIS (WOR) & 44.32\% & 0.0\% & 1,000 & 100\% \\
 & \textbf{ACTIS IS (Principled)} & \textbf{44.32\%} & \textbf{0.0\%} & \textbf{1,000} & \textbf{100\%} \\
\midrule
ImageNet & BARGAIN-PR & 100.00\% & 0.0\% & 3,000 & 100\% \\
($N=50,000$) & ACTIS (WOR) & 100.00\% & 0.0\% & 1,000 & 95\% \\
 & \textbf{ACTIS IS (Principled)} & \textbf{100.00\%} & \textbf{0.0\%} & \textbf{1,000} & \textbf{100\%} \\
\bottomrule
\end{tabular*}
\end{table}

Key insights from Table~\ref{tab:results_delta_20}:
\begin{itemize}[leftmargin=*]
    \item \textbf{Restoring Superiority over BARGAIN-PR}: On OntoNotes, the principled diagnostic reduces the oracle rate from $47.56\%$ to $42.26\%$, successfully beating BARGAIN-PR ($42.93\%$).
    \item \textbf{Superiority and Soundness on TACRED}: On TACRED, ACTIS IS achieves a $36.57\%$ oracle rate with \textbf{$0.0\%$ failures}, whereas BARGAIN-PR suffers a \textbf{$10.0\%$ empirical failure rate} while requiring more oracle calls ($40.93\%$).
    \item \textbf{Massive Gains on Large-Scale Data (Jackson)}: On Jackson, ACTIS IS requires only $44.32\%$ oracle queries, outperforming BARGAIN-PR ($76.87\%$) by \textbf{$32.55\%$ absolute}, saving over $316,000$ oracle queries.
    \item \textbf{Safety on Extreme Sparsity (ImageNet)}: On ImageNet ($50$ positives in $50,000$), all methods recognize the impossibility of aggressive filtering, safely returning $100\%$ oracle queries with $0.0\%$ error.
\end{itemize}

\subsection{Benchmark Performance at $\delta = 0.05$}
Table~\ref{tab:results_delta_05} reports performance under stricter confidence $\delta = 0.05$.

\begin{table}[ht]
\centering
\small
\setlength{\tabcolsep}{5pt}
\caption{Empirical benchmark comparison at $\gamma_R = \gamma_P = 0.95$, $\delta = 0.05$ (20 trials).}\label{tab:results_delta_05}
\begin{tabular*}{\textwidth}{@{\extracolsep{\fill}}llrrrr}
\toprule
Dataset & Method & Oracle Calls & Fail Rate & Mean $n_{\text{calib}}$ & Clean Stop \\
\midrule
OntoNotes & BARGAIN-PR & 56.64\% & 0.0\% & 1,675 & 100\% \\
($N=11,165$) & ACTIS (WOR) & 39.49\% & 0.0\% & 1,116 & 80\% \\
 & \textbf{ACTIS IS (Principled)} & \textbf{42.26\%} & \textbf{0.0\%} & \textbf{2,830} & \textbf{100\%} \\
\midrule
TACRED & BARGAIN-PR & 49.95\% & 0.0\% & 2,263 & 100\% \\
($N=22,631$) & ACTIS (WOR) & 40.03\% & 0.0\% & 2,263 & 100\% \\
 & \textbf{ACTIS IS (Principled)} & \textbf{38.37\%} & \textbf{0.0\%} & \textbf{3,421} & \textbf{95\%} \\
\midrule
Jackson & BARGAIN-PR & 83.65\% & 0.0\% & 3,000 & 100\% \\
($N=973,085$) & ACTIS (WOR) & 44.97\% & 0.0\% & 1,000 & 100\% \\
 & \textbf{ACTIS IS (Principled)} & \textbf{45.17\%} & \textbf{0.0\%} & \textbf{1,000} & \textbf{100\%} \\
\midrule
ImageNet & BARGAIN-PR & 100.00\% & 0.0\% & 3,000 & 100\% \\
($N=50,000$) & ACTIS (WOR) & 100.00\% & 0.0\% & 1,000 & 35\% \\
 & \textbf{ACTIS IS (Principled)} & \textbf{100.00\%} & \textbf{0.0\%} & \textbf{1,000} & \textbf{100\%} \\
\bottomrule
\end{tabular*}
\end{table}

At confidence $\delta = 0.05$, ACTIS Adaptive IS with the principled diagnostic strictly outperforms BARGAIN-PR across all benchmarks (by $14.38\%$ on OntoNotes, $11.58\%$ on TACRED, and $38.48\%$ on Jackson), with zero empirical constraint violations.

\newpage
\section{Conclusion and Recommendations}

The findings of this research note resolve the theoretical and empirical questions regarding asymptotic confidence sequences in adaptive model cascades:
\begin{enumerate}[leftmargin=*]
    \item \textbf{Martingale CLT Grounding}: Sequential stopping in cascade calibration is governed by the Martingale Central Limit Theorem of Hall and Heyde (1980). Valid asymptotic inference requires checking quadratic variation accumulation, conditional Lindeberg jump dispersion, and Berry-Esseen event counts.
    \item \textbf{Kish's Formula Clarification}: Kish's Effective Sample Size is a static survey concept that does not apply to sequential martingales. The variance accumulator $V_t$ already incorporates squared importance weights.
    \item \textbf{Root Cause Resolution}: The observed inefficiency of importance sampling was driven by an artificial scale divergence between $V_t$ and $v_{0, \text{floor}}$. Comparing $V_t$ against the expected process variance $v_{0, \text{target}} = n_0 \hat{\sigma}^2$ resolves this pathology cleanly without ad-hoc branching on $n_{\text{FN}} = 0$.
    \item \textbf{Codebase Implementation}: The \path{AsymptoticValidityDiagnostics} class in \path{src/actis/tuner.py} has been updated to implement Equations~\eqref{eq:check_var_r}--\eqref{eq:check_jump_p}.
\end{enumerate}

\begin{thebibliography}{99}
\bibitem{brown1971}
B.~M. Brown.
\newblock Martingale central limit theorems.
\newblock \emph{The Annals of Mathematical Statistics}, 42(1):59--66, 1971.

\bibitem{doucet2000sequential}
A.~Doucet, S.~Godsill, and C.~Andrieu.
\newblock On sequential Monte Carlo sampling methods for Bayesian filtering.
\newblock \emph{Statistics and Computing}, 10(3):197--208, 2000.

\bibitem{hall1980martingale}
P.~Hall and C.~C. Heyde.
\newblock \emph{Martingale Limit Theory and Its Application}.
\newblock Probability and Mathematical Statistics. Academic Press, New York, 1980.

\bibitem{howard2021}
S.~R. Howard, A.~Ramdas, J.~McAuliffe, and J.~Sekhon.
\newblock Time-uniform, nonparametric, nonasymptotic confidence sequences.
\newblock \emph{The Annals of Statistics}, 49(2):1055--1080, 2021.

\bibitem{kang2020approximate}
D.~Kang, P.~Bailis, and M.~Zaharia.
\newblock Model assertions for monitoring and improving ML models.
\newblock In \emph{Proceedings of Machine Learning and Systems (MLSys)}, volume~2, pages 481--496, 2020.

\bibitem{kish1965survey}
L.~Kish.
\newblock \emph{Survey Sampling}.
\newblock John Wiley \& Sons, New York, 1965.

\bibitem{lai1976boundary}
T.~L. Lai.
\newblock Boundary crossing probabilities for sample sums and Brownian motion.
\newblock \emph{The Annals of Probability}, 4(2):299--312, 1976.

\bibitem{mcleish1974}
D.~L. McLeish.
\newblock Dependent central limit theorems and invariance principles.
\newblock \emph{The Annals of Probability}, 2(4):620--628, 1974.

\bibitem{robbins1970}
H.~Robbins.
\newblock Statistical methods related to the law of the iterated logarithm.
\newblock \emph{The Annals of Mathematical Statistics}, 41(5):1397--1409, 1970.

\bibitem{ville1939etude}
J.~Ville.
\newblock \emph{{\'E}tude critique de la notion de collectif}.
\newblock Gauthier-Villars, Paris, 1939.

\bibitem{waudby2024}
I.~Waudby-Smith and A.~Ramdas.
\newblock Estimating means of bounded random variables by betting.
\newblock \emph{Journal of the Royal Statistical Society Series B: Statistical Methodology}, 86(1):1--27, 2024.
\end{thebibliography}

\end{document}
